\documentclass[11pt,a4paper]{article}
\usepackage[margin=1.9cm]{geometry}
\usepackage{amsmath,amssymb,booktabs,enumitem}
\usepackage[expansion=false]{microtype}
\usepackage[hidelinks]{hyperref}
\usepackage{xcolor}
\linespread{0.98}
\newcommand{\code}[1]{\texttt{\small #1}}
\newcommand{\flag}[1]{\textcolor{red!70!black}{#1}}
\begin{document}
\begin{center}
{\large\bfseries ½$\langle$111$\rangle\rightarrow\langle$100$\rangle$ loop conversion:
RadCluster versus Gao, Gaganidze \& Aktaa (2022)}\\[2pt]
{\footnotesize \code{RadCluster\_2\_1} checked against JNM \textbf{559} (2022) 153409 and
Acta Mater.\ \textbf{233} (2022) 117983 --- \today}
\end{center}
\vspace{2pt}

\section{Energetics: same functional family, different parameterisation}

\code{py\_utils/loop\_energetics.py} implements the Dudarev (PRL \textbf{100}, 135503)
prismatic-loop free energy
\[
E^X_l(n,T)=P_X(n)\Bigl[\hat F_X(T)\ln\!\frac{4R^*_X(n)}{e\delta}+F^X_\delta+F^X_c\Bigr],
\qquad
\hat F_{100}(T)=\hat F^{\,0}_{100}\Bigl(1-\frac{T}{T_c}\Bigr)^{1/4},\ T_c=1185\,\mathrm{K},
\]
with $\hat F_{111}$ held \emph{temperature-independent} and the single free constant
$\hat F^{\,0}_{100}$ fixed by one calibration target, $\Delta F(n_{\rm ref},T^*)=0$.
Gao \emph{et al.} instead compute $\hat F$ and $F_\delta$ from full anisotropic
elasticity with Varshni-fitted $C_{11},C',C_{44}$ to 1100\,K, and condense the result to
\[
E_{\rm tot}(N,T)=\epsilon_0(T)\sqrt N\ln N+\epsilon_1(T)\sqrt N,\qquad
\epsilon_i(T)=b_0+b_1T+b_2T^2+b_3T^3 .
\]
Since $P_X\propto\sqrt N$ and $\ln R^*=\tfrac12\ln N+\text{const}$, \textbf{these are the same
functional form}: their Table~3 is a drop-in replacement for your digitised constants.
Two structural differences matter:

\begin{enumerate}[leftmargin=*,itemsep=1pt,topsep=2pt]
\item \flag{½$\langle$111$\rangle$ is held $T$-independent in RadCluster.} Gao's fitted
$\epsilon_0,\epsilon_1$ for ½[111] carry significant linear and cubic $T$ terms. Holding
$\hat F_{111}$ fixed attributes the \emph{entire} crossover to $\langle100\rangle$ softening.
\item RadCluster carries one ½$\langle$111$\rangle$ variant. Gao resolve
$\langle11\bar2\rangle$ and $\langle1\bar10\rangle$ side orientations, whose crossover
temperatures differ by $\sim\!80^\circ$C at 100 SIAs --- comparable to the effect being modelled.
\end{enumerate}

\paragraph{Numerical comparison.} Both treatments agree in \emph{direction}: small loops
favour $\langle100\rangle$, and the crossover temperature rises with size. The shipped
calibration, however, sits high:

\begin{center}\footnotesize\setlength{\tabcolsep}{6pt}
\begin{tabular}{lcccc}
\toprule
& $n=20$ & $n=50$ & $n=100$ & $\Delta F>0$ at 450$^\circ$C \\
\midrule
RadCluster, shipped ($T^*\!=\!450^\circ$C @ $n_{\rm ref}\!=\!50$) & $-13^\circ$C & $450^\circ$C & \flag{$556^\circ$C} & $n\le49$ \\
Recalibrated $T^*\!=\!455^\circ$C @ $n_{\rm ref}\!=\!100$ & --- & $319^\circ$C & $455^\circ$C & $n\le96$ \\
Recalibrated $T^*\!=\!375^\circ$C @ $n_{\rm ref}\!=\!100$ & --- & $215^\circ$C & $375^\circ$C & $n\le183$ \\
\midrule
Gao Fig.~3(b), 100 SIAs & --- & --- & $375$--$455^\circ$C & --- \\
\bottomrule
\end{tabular}
\end{center}

\noindent $\hat F^{\,0}_{100}=0.4799$\,eV/\AA\ as shipped. Moving the calibration target
from $(450^\circ\mathrm{C},\,n\!=\!50)$ to $(455^\circ\mathrm{C},\,n\!=\!100)$ costs one
line and brings the size-dependence onto Gao's curve; the $\langle100\rangle$-favoured
window at 450$^\circ$C roughly doubles, $n\le49\rightarrow n\le96$. The module docstring
already flags the constants as ``approximate (digitised from Dudarev Figs.~2--3)'' and the
size-dependence as a calibration decision --- this is the calibration.

\section{The kinetic channels: your two, their two}

RadCluster carries \emph{three} conversion channels (\code{py\_utils/reaction\_rates.py}):
\begin{enumerate}[leftmargin=*,itemsep=1pt,topsep=2pt]
\item $\Gamma_{\rm uni}(n)=\nu_0e^{-E_a(n)/k_BT}\bigl[1-e^{-\Delta F/k_BT}\bigr]$ --- a
\emph{direct single-loop rotation}, $E_{a0}=1.8$\,eV;
\item the Marian junction $\varphi_{\rm junc}(n,n')=\varphi_{\max}
e^{-\ln^2(n/n')/2\sigma_s^2}\,\Theta(\min(n,n')\ge n_j^{\min})$ with two-step success
$P_{\rm succ}=e^{-\Delta H_2/k_BT}/(e^{-\Delta H_2/k_BT}+e^{-\Delta H_{\rm rev}/k_BT})$,
$\Delta H_2=1.0$, $\Delta H_{\rm rev}=0.30$\,eV;
\item a \emph{size-asymmetric} absorption gate $P^{\rm abs}_{\rm succ}$ with its own lower
forward barrier ($\Delta H_2^{\rm abs}=0.36$\,eV in the workbook), justified in the code as
the case where an absorbed cluster ``inherits --- is templated by --- the host's existing
$\langle100\rangle$ character.''
\end{enumerate}

Gao \emph{et al.} test both loop--loop routes and reach an asymmetric verdict:

\begin{itemize}[leftmargin=*,itemsep=1pt,topsep=2pt]
\item \flag{The symmetric Marian junction is negligible.} ``[S]imulation results show that
this contribution to the nucleation of $\langle100\rangle$ loops is negligible\ldots due to
the low possibility of two ½$\langle$111$\rangle$ loops with an identical size and suitable
orientations to encounter each other.'' Their introduction is blunter: the mechanism
``requires the collision between two migrating ½$\langle$111$\rangle$ loops with similar
sizes and overlapping glide cylinders, making it unlikely in reality.''
\item \textbf{The asymmetric reaction dominates.} Their conclusion (5) identifies
``½$\langle$111$\rangle+\langle100\rangle$'' as ``the origin causing the SIA transfer
between loops of different types,'' with the resultant character following the
\emph{larger} loop (after Arakawa's in-situ TEM).
\end{itemize}

Your channels (2) and (3) map exactly onto their negligible and dominant routes. So your
reading that both are needed is consistent with their physics --- \textbf{but the weighting
is suspect}: $\varphi_{\rm junc}$ is Gaussian in $\ln(n/n')$, \emph{peaked at $n=n'$}, and
additionally gated by $\min(n,n')\ge n_j^{\min}$. It maximises precisely the encounter Gao
find improbable. \textbf{One-line test:} set \code{phi\_max\_junc}$\,=0$ and check whether the
$\langle100\rangle$ fraction survives on the absorption channel alone. If it does, the
junction is a tuning knob, not a mechanism; if it collapses, you have a defensible
disagreement with Gao worth stating explicitly.

\section{The unary channel is the one experiment does not support}

Gao are explicit: ``The population re-arrangement proceeds \emph{not} because a single
½$\langle$111$\rangle$ loop directly transfers its orientation into the $\langle100\rangle$
configuration. Instead, loop mutual-interactions provoke the population modification.''
Yao's heavy-ion experiments saw no direct transformation --- only the abrupt disappearance
of ½$\langle$111$\rangle$ above $465^\circ$C.

$\Gamma_{\rm uni}$ is a direct single-loop rotation. Your own code comment concedes it is
calibrated rather than derived (``calibrated to place the $f_{111}(T)$ crossover at
$\sim\!350^\circ$C at reactor dose rates''). That is defensible as a coarse-grained
effective channel standing in for unresolved processes --- but it should be labelled as
such, because a referee holding these two papers will read it as the mechanism that
experiment rules out.

\section{What Gao have that RadCluster does not}

\begin{itemize}[leftmargin=*,itemsep=1pt,topsep=2pt]
\item \flag{C15 Laves clusters --- the single largest gap.} In Gao's model
$\langle100\rangle$ loops ``nucleate mainly via the collapse of C15 clusters''; C15 also
acts as an SIA buffer that ``significantly reduces the density of $\langle100\rangle$
loops'' and modifies loop reaction rates. RadCluster has no C15 population at all, so every
$\langle100\rangle$ must come from conversion. That is a genuinely different nucleation
hypothesis, not a missing detail.
\item \textbf{Cascade overlap with pre-existing clusters} (their $F_{\rm ol}$ terms,
Eq.~9) as a transformation channel among ½$\langle$111$\rangle$, $\langle100\rangle$ and C15.
\item \textbf{Annular reaction cross-sections.} They state the spherical reaction volume
``became invalid for the 2-dimensional loops with sizes over a few nanometers'' and adopt
annular cross-sections to capture coaxial-loop strain repulsion. RadCluster uses
$A_{\rm sph}m^{1/3}$ and $A_{\rm loop}n^{1/2}$ capture-radius kernels throughout.
\item \textbf{Trap-mediated loop diffusion} (Kohnert \& Wirth), $D^{\rm dt}_l=\nu^{\rm dt}_l\lambda_l^2/2N$
with $\nu^{\rm dt}\propto$ dose rate --- an irradiation-driven random walk. RadCluster uses a
thermal $D^{\rm 1D}_n\propto n^{-s}e^{-E_m/k_BT}$ with a rotational-correlation factor.
\item \textbf{Cascade cluster spectrum:} exponential, $f^{\rm SIA}(n)=0.89e^{-0.67n}$
($\sim$54\% clustered), against your power law $C_in^{-s_i}$, $f^{\rm cl}_i=0.58$. Similar
clustered fraction, very different large-cluster tail.
\end{itemize}

Conversely, Gao exclude vacancy-type loops and report that voids barely form; RadCluster's
coupled cavity--loop system covers a regime their model does not.

\section{The strongest point of agreement}

Your experimental database encodes a \emph{dose}-driven transition
(\code{loop\_burgers\_fraction.py}: ``a/2$\langle$111$\rangle$ dom.\ ($<$4 dpa);
a$\langle$100$\rangle$ dom.\ ($>$4 dpa)''), and Gao's Fig.~14 shows the same --- the
½$\langle$111$\rangle$ SIA fraction falling steeply with dose at fixed temperature,
reaching $\sim\!0$ by 0.5 dpa at 500$^\circ$C. Both therefore say the transition is not
purely thermodynamic. Gao's conclusion (7) puts it as temperature acting \emph{through
diffusivity}: ``most ½$\langle$111$\rangle$ loops are merged with $\langle100\rangle$ loops
during their long-range movements when the temperature is over (equal) 500$^\circ$C.''
That is an argument for weighting your absorption channel over both the unary and junction
channels --- and the dose dependence is the observable that discriminates them, since a
thermodynamic crossover alone predicts none.

\paragraph{Recommended actions.}
(i) Recalibrate \code{LoopEnergetics} to $(455^\circ\mathrm{C},n_{\rm ref}=100)$ or import
Gao's Table~3 directly, and give ½$\langle$111$\rangle$ its $T$-dependence.
(ii) Run the $\varphi_{\max}\!=\!0$ sensitivity to establish whether the junction is
load-bearing.
(iii) Relabel $\Gamma_{\rm uni}$ as an effective channel, or drop it and test whether
absorption alone reproduces $f_{111}(T)$.
(iv) Decide explicitly whether the absence of C15 is a modelling choice you defend or a gap
you acknowledge --- it is the first thing these authors will look for.

\end{document}
